\subsection{线性层、MLP 与 Attention 的统一检索视角}
\phantomsection\label{intro:14-summary}
采用行向量 $x:1\times d_k$，
以及逐行堆叠的 $K:M\times d_k$、$V:M\times d_v$：

\begin{center}\small
\tagpdfsetup{table/header-rows={1}}
\begin{tabularx}{\linewidth}{@{}p{23mm}p{57mm}X@{}}
\toprule
计算 & 一次读取的行向量写法 & 关键区别\\\midrule
Linear & $x(K^\top V)=xS$ & 线性匹配，外积可相加\\
MLP & $\sigma(xK^\top)V$ & 非线性筛选，权重不必归一化\\
Attention & $\softmax(\beta xK^\top)V$ & 每条查询归一化权重，Mask 限制可见范围\\
\bottomrule
\end{tabularx}
\end{center}
缩放系数 $\beta$ 在标准点积注意力里通常取
$1/\sqrt{d_k}$。三者都可以写成加权读出的形式：
\[
\text{输出}=\sum_i \text{匹配权重}_i\,v_i.
\]
但不能由此认定它们的计算、容量和训练行为相同。

\subsubsection{Key 和 Value 究竟保存在哪里？}
从 MLP 的“记忆”视角，Key 和 Value 可由训练得到的权重表示；
从 Self-Attention 的视角，它们是当前序列经过投影产生的上下文状态。
因此 MLP 的参数记忆与 Attention 的上下文读取\textbf{来源不同}。
普通 Softmax Attention 不能直接用不带条件的固定 $S$ 精确替代。
后面的现代架构将讨论一种满足特定条件的\textbf{线性状态更新}。

\chapterchecktasks{14-summary}{比较／手算／找错／场景}{
\item \textbf{比较：}分别用一句话描述外积线性记忆、Softmax Attention 和逐 Token FFN 的匹配、读取或特征变换角色。
\item \textbf{手算：}两条 Value 是 $(1,0),(0,1)$，读取权重为 $(0.25,0.75)$，输出是什么？
\item \textbf{找错：}某人看到外积读出与 Attention 都可写成权重乘 Value，就宣布它们的权重归一化、容量与更新过程完全相同。至少指出两项差异。
\item \textbf{场景：}若模型给出一句看似事实的答案，能否只根据某 FFN 神经元被激活就断言那是一条被验证的事实？应怎样补证？
}
\chapterrecallanswer{“匹配—读取”帮助比较结构，不代表 Softmax 与线性记忆等价，更不代表模型事实已被外部验证。}


